\section{课程设计与嵌赛作品的关系}
本课程设计由本人作为队长参加的嵌赛 FPGA 创新赛道作品简化而来。两个项目都围绕 FPGA 图像采集、编码与传输展开，但面向的应用场景和系统规模不同。竞赛作品强调多摄像头接入及 FPGA、处理器之间的系统协同；课程设计则选取其中的图像处理与 MJPEG 传输思路，收缩为能够在单块教学开发板上独立运行和验证的单路系统。

\subsection{竞赛作品概述}
竞赛作品对应 26 安路赛题 3，面向云端云台多视角智能安防应用。系统以 PH1A90SEG324 FPGA 和 RK3576 处理器为主要硬件平台，接入四路 CAM005 摄像头。FPGA 端完成多路图像接入、ISP 处理和 JPEG 编码，也可按工作模式输出无损 YUYV 数据；图像经 DDR 缓存后通过 PCIe 传送至 RK3576 端的 Linux 系统，由处理器侧继续完成数据接收与系统协同。当前最新的 FPGA 工程版本为 QJS2。

与只验证单个编码模块相比，该作品需要同时处理多路输入、不同数据格式、片上存储及高速接口之间的关系。因此，其设计重点不仅是 JPEG 压缩本身，还包括多路视频链路的组织、FPGA 与处理器之间的数据传输，以及各模块联调时的状态检查。

\subsection{已完成的验证}
截至本报告整理时，竞赛作品已完成若干短时功能测试。四路 1080p JPEG 模式下，短时测试中每路帧率约为 30 fps；无损 YUYV 链路也进行了单路、双路短时测试，并结合 CRC 检查数据传输的一致性，测试帧率约为 15 fps。这些结果说明多路压缩传输和可选无损链路已取得阶段性进展，但不能据此推断所有工况均已验收。长期运行稳定性及部分扩展功能仍需继续测试和完善，上述数字仅代表相应短测条件，不与课程设计的实板连续运行统计混用。

\subsection{面向课程要求的简化}
课程设计将竞赛中的多摄像头、DDR 和 PCIe 系统缩减为单摄像头、无 DDR3 的直接传输结构：OV5640 图像经过片上行缓存、预处理和 JPEG 编码后，通过 FT232H 同步 USB 接口送至 PC，由主机程序完成协议检查和 JPEG 解码显示。这样保留了流式图像处理、跨时钟传输、帧边界控制和软硬件联调等核心问题，同时去掉多路仲裁、外部大容量帧缓存及处理器侧系统，使设计规模和验证工作与 HDL 课程实验条件相适应。

因此，本课程报告中的架构、资源和性能结论均针对简化后的 Artix-7 单路系统，而不是竞赛整机指标。竞赛项目提供了系统需求和工程实践基础，课程设计则对所选链路重新组织 RTL、测试和实板验收，以更集中地说明各模块的接口关系及其实现依据。
